\documentclass[9pt,twocolumn]{article}
\usepackage[letterpaper,top=0.58in,bottom=0.62in,left=0.62in,right=0.62in,columnsep=0.24in]{geometry}
\usepackage[T1]{fontenc}
\usepackage{newtxtext,newtxmath}
\usepackage{microtype}
\usepackage{amsmath,mathtools}
\usepackage{booktabs,tabularx,array}
\usepackage{enumitem}
\usepackage{xcolor}
\usepackage{titlesec}
\usepackage{graphicx}
\usepackage{tikz}
\usepackage{hyperref}
\usepackage{caption}
\usepackage{fancyhdr}
\usepackage{balance}
\usetikzlibrary{positioning,arrows.meta,fit}

\definecolor{ink}{RGB}{22,22,22}
\definecolor{shade}{RGB}{246,246,246}
\definecolor{rulegray}{RGB}{95,95,95}
\hypersetup{colorlinks=true,linkcolor=ink,citecolor=ink,urlcolor=ink}
\setlength{\parindent}{1em}
\setlength{\parskip}{0pt}
\setlist{nosep,leftmargin=1.35em}
\captionsetup{font=footnotesize,labelfont=bf}
\titleformat{\section}{\normalfont\bfseries\scshape\small}{\Roman{section}.}{0.45em}{}
\titleformat{\subsection}{\normalfont\bfseries\small}{\Alph{subsection}.}{0.4em}{}
\titlespacing*{\section}{0pt}{1.05ex plus .2ex minus .1ex}{0.45ex}
\titlespacing*{\subsection}{0pt}{0.8ex plus .15ex minus .1ex}{0.3ex}
\pagestyle{fancy}
\fancyhf{}
\fancyhead[L]{\scriptsize Zot \;|\; Systemic Continuity Theory of Consciousness}
\fancyhead[R]{\scriptsize The Not-So-Hard Problem of Consciousness}
\fancyfoot[C]{\scriptsize \thepage}
\renewcommand{\headrulewidth}{0.25pt}
\setlength{\headheight}{12pt}

\newcommand{\Csys}{C_B(t)}
\newcommand{\Lsys}{L_B(t)}
\newcommand{\Gam}{\Gamma_B(t_0,t_1)}
\newcommand{\Fvec}{\mathbf{F}_B(t)}
\newcommand{\Ephen}{E_B(t)}
\newcommand{\repoURL}{https://github.com/mikecreation/ZotBot/tree/main/research/not-so-hard-problem-of-consciousness}
\newcommand{\defbox}[1]{%
\vspace{0.4ex}\noindent\colorbox{shade}{\parbox{0.965\columnwidth}{\small #1}}\vspace{0.4ex}}

\begin{document}
\twocolumn[
\begin{@twocolumnfalse}
\begin{center}
{\LARGE\bfseries The Not-So-Hard Problem of Consciousness\par}
\vspace{0.25em}
{\large\itshape Defining the System Before Explaining Experience\par}
\vspace{0.65em}
{\normalsize\bfseries Michael Zot$^{1,*}$\par}
\vspace{0.15em}
{\small $^1$Independent Researcher, Illinois, USA \quad $^*$mike@ZotBot.Ai\par}
\vspace{0.12em}
{\small\href{\repoURL}{github.com/mikecreation/ZotBot/.../not-so-hard-problem-of-consciousness}\par}
\end{center}
\vspace{0.45em}
\hrule
\vspace{0.55em}
\noindent\textbf{Abstract.}
Science cannot explain a variable whose referent is allowed to change between experiments. Yet \emph{consciousness} is routinely used for organismal existence, wakefulness, responsiveness, environmental connectedness, phenomenal experience, memory, cognitive access, self-awareness, reportability, intelligence, and momentary mental content. This paper fixes the target before asking for its mechanism. The Systemic Continuity Theory of Consciousness (SCTC) defines minimal system-level consciousness as the ongoing viable existence of an integrated biological individual. Phenomenal experience is represented separately as $E$. This separation is not cosmetic. It blocks a recurrent target-substitution error: failure to explain why a dream or sensation is experienced does not count as failure to define or explain system-level consciousness unless an independent argument first establishes that consciousness and phenomenal experience are the same variable. The familiar ``hard problem of consciousness'' therefore contains a prior definitional commitment. Chalmers' hard target is experience; calling that target \emph{consciousness} already selects $C:=E$. SCTC rejects that identity and treats the hard problem, if it remains, as a hard problem of phenomenal experience. The theory is definition-first: dissociations do not prove the definition; they reveal its consequences. Anesthesia, dreaming, amnesia, sensory loss, cognitive motor dissociation, replacement, resuscitation, and external support demonstrate that many variables commonly packed into \emph{consciousness} can change while the biological individual continues. A direct rival comparison, a closed-inventory system-level discontinuity test, and explicit failure conditions make the proposal empirically vulnerable. The claim is narrow but consequential: before asking how consciousness is produced, science must first decide what it is talking about.

\vspace{0.35em}
\noindent\textbf{Keywords:} consciousness; phenomenal experience; biological continuity; anesthesia; definition; hard problem; cognitive motor dissociation; biological individuality
\vspace{0.55em}
\hrule
\vspace{0.8em}
\end{@twocolumnfalse}
]

\section{The Problem Precedes the Mechanism}
The canonical question, ``How does the brain produce consciousness?'', appears mechanistic. It is actually downstream of a more basic question: \emph{what variable does the word consciousness denote?} Chalmers himself noted that ``consciousness'' is ambiguous before isolating conscious experience as the hard target \cite{Chalmers1995}. Block later called consciousness a ``mongrel concept'' and separated phenomenal from access consciousness \cite{Block1995}. Sanders et al. defined consciousness as subjective experience while distinguishing it from responsiveness and environmental connectedness \cite{Sanders2012}. Bayne, Hohwy, and Owen rejected a single levels-based description in favor of a multidimensional account \cite{Bayne2016}. LeDoux summarized the difficulty directly: what functions are assigned to consciousness depends on what one takes consciousness to be \cite{LeDoux2025}.

These are not merely verbal disagreements. A study of responsiveness, a study of dream experience, a study of memory, and a study of organismal viability can each be rigorous while targeting different variables. If all four are reported as measurements of one thing, the noun becomes part of the experimental design.

\defbox{\textbf{Target rule.} A mechanism can be evaluated only after its dependent variable is fixed. If the referent changes, apparent explanatory failure can be produced by target substitution rather than by the phenomenon.}

SCTC therefore makes an explicit revisionary assignment rather than allowing the noun to move between experiments.

\section{Definition Lock: What This Paper Means by Consciousness}
\subsection{The biological individual and its continuity}
Let $B$ denote a biological individual. Let
\begin{equation}
\Gamma_B(t_0,t_1)=1
\end{equation}
mean that the individual at $t_1$ lies on the continuing organism-level causal trajectory of the individual at $t_0$. Material turnover, repair, component replacement, architectural reorganization, and external support are compatible with continuity when they preserve the same biological individual.

Continuity across times and viability at a time are different objects. Define the living-system state
\begin{equation}
L_B(t)=
\begin{cases}
1, & B \text{ exists as a viable integrated biological individual at }t,\\
0, & B \text{ has reached terminal organism-level discontinuity.}
\end{cases}
\end{equation}
SCTC then fixes the minimal system-level target by definition:
\begin{equation}
\boxed{C_B(t)\equiv L_B(t).}
\label{eq:def}
\end{equation}
Equation \ref{eq:def} is primitive within the theory. Later cases do not prove it. The logical direction is
\begin{equation}
\text{definition}\;\rightarrow\;\text{deductions}\;\rightarrow\;\text{comparative tests}.
\end{equation}

\subsection{Faculties are not aliases for the system}
Let the declared faculty/state vector be
\begin{equation}
\mathbf F_B^*(t)=(W,E,K,R,M,P,Q,S,I),
\end{equation}
where $W$ is wakefulness/arousal; $E$ phenomenal experience; $K$ cognitive or experiential connectedness to the environment; $R$ behavioral responsiveness; $M$ memory; $P$ perceptual/sensory capacities; $Q$ cognitive access; $S$ self-awareness/metacognition; and $I$ intelligence or related cognitive capacities. Let $\mathcal A_B(t)$ denote the architecture supporting these capacities and $X_B(t)$ the state currently instantiated.

The hierarchy is descriptive, not set-theoretic:
\begin{equation}
B\text{ continuing}\;\longrightarrow\;\mathcal A_B\;\longrightarrow\;\mathbf F_B^*\;\longrightarrow\;X_B.
\end{equation}

\defbox{\textbf{Definition lock.} Throughout this paper, $C$ means system-level consciousness as defined in Eq. \ref{eq:def}. $E$ means phenomenal experience. No argument in this paper uses $C$ as a synonym for $E$. An objection that silently substitutes $E$ for $C$ is not an internal objection to SCTC; it is a rival target assignment and must be evaluated as such.}

A resolution rule is also required. $R$ and $K$ in $\mathbf F_B^*$ are behavioral/cognitive variables. Minimal biological responsiveness or metabolic coupling used to assess viability belongs to the implementation of $L_B$, not to the behavioral variables $R$ and $K$. Identical words at different biological scales do not denote identical variables.

\section{The Target-Substitution Error}
The most persistent misunderstanding can be stated formally.

\textbf{Principle 1 (Faculty exclusion).} If a declared faculty $F_j$ can be absent while $L_B(t)=1$, then $F_j$ is not constitutive of $C_B$ under SCTC:
\begin{equation}
L_B(t)=1\;\land\;F_j(t)=0 \;\Rightarrow\; F_j\not\equiv C_B.
\end{equation}
This is a consequence of the definition, not evidence that generated it.

\textbf{Principle 2 (Explanatory separation).} For any nonconstitutive faculty $F_j$,
\begin{equation}
F_j\not\equiv C_B
\quad\Rightarrow\quad
\mathrm{Unexplained}(F_j)\not\Rightarrow\mathrm{Unexplained}(C_B).
\label{eq:sep}
\end{equation}
Equation \ref{eq:sep} is the firewall between the theory and a recurring category error. If experience is not constitutive of system-level consciousness, then failing to explain why a particular physical process was experienced cannot by itself refute the definition of consciousness.

\subsection{The anesthesia dream, without the fog}
Consider an anesthetized person whose organismal continuity is preserved. Three epistemic cases are possible:
\begin{align}
L_B=1,\;R=0,\;E=1 &\quad \text{experience occurred},\\
L_B=1,\;R=0,\;E=0 &\quad \text{no experience occurred},\\
L_B=1,\;R=0,\;E=? &\quad \text{experience is unknown}.
\end{align}
Under SCTC, all three satisfy $C_B=1$. The scientifically valid question ``Why did a dream occur in the first case?'' targets $E$. It may be extremely difficult. It does not become a counterexample to Eq. \ref{eq:def} merely because the historical vocabulary called $E$ ``consciousness.''

This distinction can be summarized in one sentence:

\defbox{\textbf{A hard question about a faculty does not become a hard question about the bearer merely because both were given the same noun.}}

\section{Empirical Dissociations: What the Data Actually Separate}
The dissociation literature does not establish Eq. \ref{eq:def}; it demonstrates why the variables packed into the conventional noun should not be treated as interchangeable.

\subsection{Responsiveness, connectedness, and experience}
Sanders et al. explicitly argued that unresponsiveness is not equivalent to unconsciousness and distinguished subjective experience from responsiveness and connectedness \cite{Sanders2012}. Noreika et al. induced unresponsiveness with sedative/anesthetic agents and found later reports of subjective experience in almost 60\% of sessions \cite{Noreika2011}. Casey et al. evaluated EEG measures against explicit definitions of responsiveness, connectedness, and subjective experience, finding that measures derived from unresponsiveness did not map uniquely onto consciousness alone \cite{Casey2024}. Therefore
\begin{equation}
R=0\not\Rightarrow E=0.
\end{equation}
SCTC adds only the system-level statement that ordinary reversible anesthesia is compatible with $L_B=1$ while the faculty vector changes.

\subsection{Observable response and cognitive capacity}
Bodien et al. studied 353 adults with disorders of consciousness in a prospective, multicenter cohort. Among 241 participants without an observable response to commands, 60 (25\%) showed task-based command following on fMRI or EEG \cite{Bodien2024}. Thus
\begin{equation}
R_{\mathrm{observable}}=0\not\Rightarrow Q_{\mathrm{command}}=0.
\end{equation}
The scanner did not create the capacity. It exposed a capacity that the behavioral channel had failed to reveal.

\subsection{Memory and sensation}
Severe amnesia can coexist with preserved perception, language, reasoning, and continued biological existence \cite{Milner2005}. Individual sensory modalities can likewise be lost while the organism persists. These are elementary but decisive separations: memory, vision, hearing, or reportability may characterize the state of a conscious system without constituting the system-level category itself.

\begin{table*}[t]
\centering
\caption{The same continuing individual can occupy radically different faculty states. The table is not evidence for the primitive definition; it shows why the variables should not be used as synonyms.}
\footnotesize
\begin{tabularx}{0.96\textwidth}{@{}l c c c c X@{}}
\toprule
Condition & $L_B$ & $R$ & $E$ & $M$ & What is separated \\
\midrule
Ordinary wakefulness & 1 & often 1 & often 1 & variable & baseline faculty configuration \\
Reversible anesthesia & 1 & 0 & 0/1/? & often impaired & system continuity from response, experience, recall \\
Dreaming sleep & 1 & low & often 1 & variable & experience from environmental response \\
Severe amnesia & 1 & variable & possible & severely impaired & memory from system continuity \\
Sensory loss & 1 & variable & modality-specific loss & variable & modality from system continuity \\
Cognitive motor dissociation & 1 & 0 by behavior & unknown/possible & variable & observable response from cognitive command following \\
\bottomrule
\end{tabularx}
\end{table*}

\section{Why the ``Hard Problem of Consciousness'' Is Not a Neutral Phrase}
Chalmers' hard problem is the problem of why physical processing is accompanied by conscious experience \cite{Chalmers1995}. That is a coherent target. Under the notation here, it is the problem of $E$.

The phrase \emph{hard problem of consciousness} is therefore not definition-free. It becomes neutral only if $C\equiv E$ has already been established. Otherwise the phrase imports a contested identity into the name of the problem.

Under SCTC the correct question is
\begin{equation}
\boxed{\text{Which physical organizations instantiate }E,\text{ and why?}}
\end{equation}
That may remain the deepest problem in the science of mind. SCTC does not answer it by renaming it. SCTC denies that its unresolved status licenses the inference that system-level consciousness is likewise undefined.

The distinction is exactly the one exposed by the anesthesia example. Asking why a dream was experienced is a question about $E$. Asking whether the biological individual continued is a question about $L_B$. They are not two phrasings of the same dependent variable.


\section{Support Is Not Identity}
A component can be necessary for maintaining a system without being identical to the system-level property.

\textbf{Proposition 1 (Support is not identity).} If component $x$ is causally necessary for maintaining $B$ under a specified condition, it does not follow that $x\equiv B$ or $x\equiv C_B$.

A heart can be necessary for circulation without being the organism. Particular neural structures can be necessary for wakefulness, breathing, memory, or reportability without being identical to the continuing individual. Organ replacement makes the distinction visible: if a component is replaced while organism-level continuity is preserved, the individual has not been recreated merely because the implementation changed.

Temporary functional collapse is likewise not identical to terminal discontinuity. Let $\mathcal I$ be a declared family of admissible rescue interventions. Define the terminal boundary
\begin{equation}
\tau_B^*=\inf\left\{t:\substack{\text{no admissible trajectory under }\mathcal I\\
\text{preserves }B\text{ as the same viable individual}}\right\}.
\end{equation}
The criterion is explicitly counterfactual and intervention-relative. Successful resuscitation is evidence that the terminal boundary had not been crossed under the declared intervention set; it is not a proof that all imaginable identity questions are settled.

\section{Biological Individuality Is a Dependency, Not a Hidden Escape}
SCTC inherits rather than solves the biological-individuality problem. Fulda treats strong agential autonomy as a criterion of biological individuality \cite{Fulda2023}; Bechtel and Bich analyze physiological regulation from the perspective of maintaining the organism through time \cite{BechtelBich2024}. These frameworks motivate, but do not completely settle, every boundary case.

Accordingly, where biology can identify an integrated, relatively autonomous individual with acceptable intersubjective reliability, SCTC inherits that classification. Where individuality is genuinely indeterminate, as may occur in some cases involving fission, fusion, pregnancy, colonial organization, cryptobiosis, or intensive machine support, SCTC inherits the indeterminacy. It does not replace an unresolved biological boundary with a neurological stipulation.

This matters for brain death. The 2023 AAN/AAP/CNS/SCCM guideline provides clinical criteria for death by neurologic criteria \cite{Greer2023}. Shewmon documented prolonged somatic survival after formal brain-death diagnoses and argued that absence of brain function does not by itself establish loss of all somatic integrative unity \cite{Shewmon1998}. SCTC is not medical or legal guidance. Its theoretical commitment is narrower:

\textbf{Consequence 1.} If a permanently brain-nonfunctional body still satisfies the best operational criteria for one viable integrated biological individual, then $L_B=1$ and therefore $C_B=1$ under SCTC. If the remaining physiology no longer constitutes one biological individual, then $L_B=0$ and $C_B=0$. If biological individuality is indeterminate, SCTC is indeterminate.

\section{Rival Assignments Must Be Named, Not Smuggled In}
SCTC is not the only coherent assignment.

\subsection{Experience-centered rival}
The nearest naming rival is
\begin{equation}
\mathcal R_E:\qquad C:=E,
\end{equation}
with responsiveness, connectedness, memory, access, reportability, and other variables decomposed. Sanders et al. instantiate this style explicitly \cite{Sanders2012}. Shared success on anesthesia or dissociation does not adjudicate SCTC, because both assignments predict those separations.

\subsection{Capacity-transition rivals}
A biologically substantive rival places the transition to minimal consciousness at a specific capacity rather than at organismal viability. Unlimited Associative Learning (UAL), for example, has been proposed as a transition marker for minimal consciousness or sentience \cite{Bronfman2016,Birch2020}. SCTC does not deny that such capacities may mark the emergence of sentience or phenomenal experience. It denies that a capacity marker is thereby the persistence condition of the biological individual.

The disagreement is therefore explicit:
\begin{align}
\mathcal R_\Gamma &: C:=L_B,\\
\mathcal R_E &: C:=E,\\
\mathcal R_T &: C:=T \quad\text{for a proposed capacity transition }T.
\end{align}
No rival may be used as an unstated definition while simultaneously being offered as an objection.


\section{Comparative Tests and Failure Conditions}
\subsection{Three-arm target-assignment experiment}
Randomize clinicians, researchers, or trained participants to identical longitudinal cases under three vocabularies: umbrella language (Arm U), experience-centered decomposition (Arm E), and SCTC decomposition (Arm $\Gamma$). All arms must answer the same forced response fields:
\begin{enumerate}
\item individual status: same individual / ended / indeterminate;
\item phenomenal experience: present / absent / unknown;
\item responsiveness: present / absent / unknown;
\item connectedness: present / absent / unknown.
\end{enumerate}
The missingness mechanism must be preregistered. ``Unknown'' is not scored as a reclassification.

Let $V_R$ be the number of unsupported flips in the system-level classification caused solely by changes in evidence availability, and let $\mathrm{Err}_E(R)$ be error on separately scored phenomenal-state judgments. The distinctive SCTC prediction is conjunctive:
\begin{equation}
\boxed{V_\Gamma<V_E\quad\land\quad \mathrm{Err}_E(\Gamma)-\mathrm{Err}_E(E)\le\varepsilon_E.}
\end{equation}
Reduced volatility is not a victory if it is purchased by making people worse at reasoning about experience.

\subsection{Closed-inventory system-level discontinuity test}
To prevent post hoc absorption of counterexamples, the faculty inventory is frozen in advance as $\mathbf F_B^*=(W,E,K,R,M,P,Q,S,I)$. Search across at least two preregistered paradigms for a variable $Z$ such that:
\begin{enumerate}
\item $Z$ changes while $L_B=1$;
\item $Z$ is not reducible to $\mathbf F_B^*$ or measurement availability;
\item $Z$ adds independent predictive or interventional information after those variables are controlled;
\item the relevant outcome family and reduction criteria were fixed before observing $Z$.
\end{enumerate}
If such a system-level $Z$ is necessary to recover the classifications science intends by consciousness, then $C\equiv L_B$ is insufficient and SCTC fails at its core.

\subsection{Explicit failure conditions}
SCTC should be abandoned or revised if: (i) the direct assignment experiment favors $\mathcal R_E$ or another rival; (ii) a preregistered independent system-level $Z$ survives the discontinuity test; (iii) biological individuality cannot be operationalized with adequate reliability for the intended domain; or (iv) the biological substrate restriction fails because a nonbiological system satisfies the deeper organizational criterion that the theory actually requires.

\section{What SCTC Does Not Entail}
Because the paper uses \emph{consciousness} at the minimal system level, several non-entailments must be explicit:
\begin{equation}
C_B=1\;\not\!\Rightarrow\;\substack{E=1,\;\text{pain, sentience, interests, personhood,}\\
\text{moral status, or legal status}.}
\end{equation}
Those are separate empirical or normative variables. The point is not to weaken them. It is to prevent them from being re-imported into $C$ after the theory has separated them.

The same rule applies to clinical language. ``Loss of consciousness'' should be replaced by the variable actually measured: loss of responsiveness, wakefulness, connectedness, memory formation, reportability, or evidence of phenomenal experience. ``Neural correlate of consciousness'' should identify the particular faculty or state whose correlate is being measured.

\section{Discussion: The Problem Was Allowed to Move}
The strongest criticism of SCTC is that experience was the target researchers cared about all along. That may often be true. It does not establish that experience and system-level consciousness are the same variable. It establishes that experience deserves its own theory.

This is the central methodological claim. A question can remain profound after a terminology correction. What changes is which variables inherit the mystery. If the unresolved question is why physical organization gives rise to subjective experience, then $E$ is the target. Study $E$ as $E$. Do not require wakefulness, memory, reportability, responsiveness, sensory access, or organismal continuity to inherit the same explanatory debt because one historical noun covered them all.

The paper therefore does not claim to have explained dreaming, pain, visual phenomenology, or the existence of a point of view. It claims that none of those unanswered questions can be used against Eq. \ref{eq:def} without first supplying an argument that the faculty in question is identical to the system-level variable.

\defbox{\textbf{The decisive question is not ``Have you explained experience?'' It is ``Why should an unexplained theory of experience count as an unexplained definition of the continuing system?''}}

That question cannot be answered by changing the meaning of consciousness halfway through the argument.

\section{Conclusion}
A scientific mechanism requires a stable target. Consciousness research inherited a word that has been used for a biological individual, its wakefulness, its responses, its access to the environment, its memories, its self-model, its reports, and its phenomenal experiences. Those variables do not move together.

SCTC fixes the system-level target first:
\begin{equation}
\boxed{\begin{aligned}
C_B(t)&\equiv L_B(t)\\
&=\text{viable organism-level existence of }B.
\end{aligned}}
\end{equation}
Everything else must earn its own variable.

This does not solve the hard problem of phenomenal experience. It prevents that problem from being silently substituted for every other question called consciousness. The anesthesia dream is the simplest demonstration. The person continues whether a dream occurs, does not occur, or cannot later be recovered. Why a dream occurred is a real scientific question about $E$. It is not evidence that the continuing individual disappeared and returned.

The historical hard problem may therefore contain two different things: a genuinely difficult problem of phenomenal experience and a prior representation error about what object the word \emph{consciousness} was supposed to denote. The first deserves explanation. The second deserves correction.

\begin{center}
\small\bfseries A mechanism can be mysterious only after its target is stable.
\end{center}


\balance
\begin{thebibliography}{99}\small
\bibitem{Chalmers1995} D. J. Chalmers, ``Facing up to the problem of consciousness,'' \emph{J. Conscious. Stud.}, vol. 2, no. 3, pp. 200--219, 1995.
\bibitem{Rosenthal1993} D. M. Rosenthal, ``State consciousness and transitive consciousness,'' \emph{Conscious. Cogn.}, vol. 2, no. 4, pp. 355--363, 1993, doi:10.1006/ccog.1993.1029.
\bibitem{Block1995} N. Block, ``On a confusion about a function of consciousness,'' \emph{Behav. Brain Sci.}, vol. 18, no. 2, pp. 227--247, 1995, doi:10.1017/S0140525X00038188.
\bibitem{Sanders2012} R. D. Sanders, G. Tononi, S. Laureys, and J. W. Sleigh, ``Unresponsiveness $\neq$ unconsciousness,'' \emph{Anesthesiology}, vol. 116, no. 4, pp. 946--959, 2012, doi:10.1097/ALN.0b013e318249d0a7.
\bibitem{Bayne2016} T. Bayne, J. Hohwy, and A. M. Owen, ``Are there levels of consciousness?'' \emph{Trends Cogn. Sci.}, vol. 20, no. 6, pp. 405--413, 2016, doi:10.1016/j.tics.2016.03.009.
\bibitem{LeDoux2025} J. E. LeDoux, ``What the functions of consciousness are depends on what one thinks consciousness is,'' \emph{Philos. Trans. R. Soc. B}, vol. 380, art. 20240311, 2025, doi:10.1098/rstb.2024.0311.
\bibitem{Noreika2011} V. Noreika et al., ``Consciousness lost and found: subjective experiences in an unresponsive state,'' \emph{Brain Cogn.}, vol. 77, no. 3, pp. 327--334, 2011, doi:10.1016/j.bandc.2011.09.002.
\bibitem{Casey2024} C. P. Casey et al., ``Evaluation of putative signatures of consciousness using specific definitions of responsiveness, connectedness, and consciousness,'' \emph{Br. J. Anaesth.}, vol. 132, no. 2, pp. 300--311, 2024, doi:10.1016/j.bja.2023.09.031.
\bibitem{Bodien2024} Y. G. Bodien et al., ``Cognitive motor dissociation in disorders of consciousness,'' \emph{N. Engl. J. Med.}, vol. 391, pp. 598--608, 2024, doi:10.1056/NEJMoa2400645.
\bibitem{Milner2005} B. Milner, ``The medial temporal-lobe amnesic syndrome,'' \emph{Psychiatr. Clin. North Am.}, vol. 28, no. 3, pp. 599--611, 2005, doi:10.1016/j.psc.2005.06.002.
\bibitem{Fulda2023} F. C. Fulda, ``Agential autonomy and biological individuality,'' \emph{Evol. Dev.}, vol. 25, no. 6, pp. 353--370, 2023, doi:10.1111/ede.12450.
\bibitem{BechtelBich2024} W. Bechtel and L. Bich, ``Situating homeostasis in organisms: maintaining organization through time,'' \emph{J. Physiol.}, vol. 602, no. 22, pp. 6003--6020, 2024, doi:10.1113/JP286883.
\bibitem{Bronfman2016} Z. Z. Bronfman, S. Ginsburg, and E. Jablonka, ``The transition to minimal consciousness through the evolution of associative learning,'' \emph{Front. Psychol.}, vol. 7, art. 1954, 2016, doi:10.3389/fpsyg.2016.01954.
\bibitem{Birch2020} J. Birch, S. Ginsburg, and E. Jablonka, ``Unlimited associative learning and the origins of consciousness: a primer and some predictions,'' \emph{Biol. Philos.}, vol. 35, art. 56, 2020, doi:10.1007/s10539-020-09772-0.
\bibitem{Greer2023} D. M. Greer et al., ``Pediatric and adult brain death/death by neurologic criteria consensus guideline,'' \emph{Neurology}, vol. 101, no. 24, pp. 1112--1132, 2023, doi:10.1212/WNL.0000000000207740.
\bibitem{Shewmon1998} D. A. Shewmon, ``Chronic `brain death': meta-analysis and conceptual consequences,'' \emph{Neurology}, vol. 51, no. 6, pp. 1538--1545, 1998, doi:10.1212/WNL.51.6.1538.
\end{thebibliography}

\vspace{0.5em}
\noindent\footnotesize\textbf{Methodological note.} This article is a theoretical and methodological proposal. It reports no new human or animal experiment and is not clinical, legal, or ethical guidance. Empirical claims are limited to the cited literature. The definition in Eq. \ref{eq:def} is revisionary and is evaluated by coherence, comparative scientific utility, and the explicit failure conditions stated above; it is not presented as a result established by the dissociation cases themselves.\\[0.35em]\textbf{Open materials.} Manuscript PDF, LaTeX source, experiment specifications, reviewer-response notes, social-post copy, and explainer-video materials are archived at \href{\repoURL}{the project repository}.

\end{document}

